% FILE: 03_architecture_and_primitives.tex
% Project: research-open-ontologies
% Role: open-ontology-research-and-rdf-authority

\section{Architecture and Primitives}
\label{sec:research-open-ontologies-architecture}

\subsection{Core Objects}
\label{sec:research-open-ontologies-core-objects}

The ontology corpus is organized into several layers of named objects, each grounded in a
specific TTL file validated by the Phase~7 roundtrip report.

\paragraph{M\&A Claim Hierarchy (\texttt{ma:BoardClaim}).}
Defined as an OWL class hierarchy with five subclasses: \texttt{SynergyProjection},
\texttt{OperationalDebtClaim}, \texttt{IntegrationRiskClaim}, \texttt{ProcessAssetClaim},
and \texttt{ControlClaim}. Each subclass requires a named receipt type and a named witness
marker before the claim is board-admissible. This hierarchy is the formal type system for
the 42 M\&A claim files in the \texttt{ma/} directory.

\paragraph{Conformance Verdict (\texttt{wasm4pm:ConformanceVerdict}).}
A typed individual requiring \texttt{fitness >= 0.95} and \texttt{precision >= 0.90},
backed by a \texttt{BLAKE3 CryptographicReceipt}. This is the minimum threshold for a
process intelligence claim to be presented at board level.

\paragraph{Witness Lattice (\texttt{compat:WitnessMarker}).}
Eight named witnesses defined in \path{ggen/wasm4pm-compat.ttl}, forming a join-semilattice:
VanDerAalst1989 (Petri net marking, lattice position \texttt{bottom}), VanDerAalst1998
(soundness proof), VanDerAalst2016 (OCEL object-centric), Murata1989 (Petri net semantics,
\texttt{bottom}), Weijters2011 (causal net / Heuristics Miner), Leemans2013 (process tree /
Inductive Miner), Adriansyah2011 (alignment-based conformance), and BlueRiverDam (autonomic
actuation, lattice position \texttt{top}). Every conformance claim must name a witness from
this lattice.

\paragraph{Type-Law Atoms (\texttt{compat:TypeLawAtom}).}
Six named laws: 1BoundednessLaw, TerminalMarkingLaw, ObjectInstanceLaw, FitnessThresholdLaw,
RoleBasedSignatureLaw, and JcsCanonicalLaw. These laws are the admission predicates for
process evidence at the compatibility layer boundary.

\paragraph{Boundary Crossings (\texttt{compat:BoundaryCrossing}).}
Nine named crossings: LogToPetriNet, LogToProcessTree, OcelToObjectCentric,
ConstraintEvaluation, TokenReplay, Alignment, ReceiptGraduation, AutonomicActuation,
WitWorldExport. Each crossing defines a typed boundary in the execution architecture.

\paragraph{Lifecycle States (\texttt{lifecycle:ProcessState}).}
Seven states: Design, Simulation, Validation, Monitoring, Optimization, Repair, Decommission.
Governed by MAPE-K rules defined in \path{autonomic-law.ttl}.

\paragraph{Program Role Taxonomy (\texttt{pi:ProgramRole}).}
Eight roles defined in \path{pi-program.ttl}: PROGRAM, PROOF\_CELL, ENGINE,
COMPATIBILITY\_LAYER, MANUFACTURING\_CELL, TELEMETRY\_FEEDSTOCK, MOBILE\_SUBSTRATE,
WORKFLOW\_SUBSTRATE. Every artifact in the research program is classified under one of these
roles.

\subsection{Admissible Transitions}
\label{sec:research-open-ontologies-transitions}

The MAPE-K loop defined in \path{autonomic-law.ttl} governs all transitions between
lifecycle states. The transition partition is binary:

\begin{itemize}
  \item \texttt{aka:ElasticTransition} --- autonomous local repair and scaling operations
        that do not require a governance token (\texttt{aka:requiresGovToken false}).
        Subtypes: LocalRepair (single S-component modification maintaining boundary
        preservation), ScalingTuning (queue capacity, resource scaling, logging adjustments).
  \item \texttt{aka:ComplianceTransition} --- policy mutation and structural replacement
        operations that require a governance token (\texttt{aka:requiresGovToken true}).
        Subtypes: PolicyMutation (modify global LTL safety policies $\Phi_{Gov}$),
        StructuralReplacement (global model hot-swap $W \to W'$), LifecyclePromotion
        (move models to active status or decommission).
\end{itemize}

The primary containment invariant is the LTL safety invariant:
\[
  \texttt{gov:LTLSafetyInvariant:}\ G(\neg\mathit{compliant} \implies \mathbf{X}(\neg\mathit{actuated}))
\]
This invariant is encoded in \path{governance-policy.ttl} and states: globally, if a process
state is not compliant, then at the next step it must not be actuated. No autonomic action
may actuate a non-compliant process.

\subsection{Boundaries}
\label{sec:research-open-ontologies-boundaries}

The ontology corpus enforces two categories of boundaries:

\paragraph{DTO/JSON Boundary (\texttt{pi:FORBIDDEN\_COLLAPSE}).}
Defined as a law: no data transfer object or JSON serialization method may cross the
compatibility layer boundary. JSON serialization belongs inside the wasm4pm ENGINE, not in
the COMPATIBILITY\_LAYER. The conformance audit (PI\_GGEN\_UNIFIED\_RUN\_CONFORMANCE\_AUDIT\_001)
identified two violations of this law (\texttt{to\_json\_string()} and \texttt{receipt\_json()}
in \path{compat/src/manufacturing/}), confirmed as Gate~1 FAIL in the GGEN conformance audit.

\paragraph{Namespace Boundary.}
The ontology files use two namespace regimes: local extension namespaces
(\texttt{http://process-intelligence.local/} for autonomic law, MAPE-K) and
published-prefix namespaces (\texttt{https://process.intelligence/} for the pi-program,
compat, wasm4pm, lifecycle, and ma ontologies). The roundtrip report exposed a SPARQL
prefix mismatch: \texttt{count\_checkpoints=0} despite 3{,}684 triples loaded, indicating
that ontology entities are present in the graph but the SPARQL queries used the wrong
prefix to address them. This is a known open gap.

\subsection{Failure Modes Refused}
\label{sec:research-open-ontologies-failure-modes}

The following failure modes are explicitly refused by the architecture:

\begin{description}
  \item[Circular dependency in TTL imports.] The federated dependency graph enforces
        \texttt{ACYCLIC} status as a gate criterion. A circular import chain would cause
        the gate to fail, blocking all downstream manufacturing. As of 2026-06-01, zero
        circular dependencies exist.
  \item[Dangling namespace references.] The dependency graph validates that every namespace
        prefix used in a TTL file resolves to a loaded ontology. Zero dangling references
        were found in the Phase~4 analysis.
  \item[Forced ALIVE on partial evidence.] The program COVENANT explicitly forbids declaring
        ALIVE without all gate criteria met. The open-ontologies project is correctly
        classified as PARTIAL pending TTL syntax validation via \texttt{rapper} and receipt
        TTL manufacture.
  \item[Ad hoc conformance claims.] A claim of fitness or precision without a named
        \texttt{compat:WitnessMarker} and a BLAKE3 receipt is inadmissible. The ontology
        corpus provides the type system that makes this refusal enforceable at the RDF layer.
\end{description}
